\documentclass[acmsmall,nonacm,screen]{acmart}

\usepackage{mathpartir}
\usepackage{booktabs}

\citestyle{acmauthoryear}
\settopmatter{printacmref=false}

\AtEndPreamble{\theoremstyle{acmdefinition}\newtheorem{remark}[theorem]{Remark}}

\newcommand{\kw}[1]{\mathsf{#1}}
\newcommand{\nm}[1]{\mathit{#1}}
\newcommand{\Base}{\kw{Base}}
\newcommand{\Ok}[2]{\kw{Ok}_{#1}(#2)}
\newcommand{\Bud}[2]{\kw{Bud}_{#1}(#2)}
\newcommand{\ok}[2]{\kw{ok}_{#1}(#2)}
\newcommand{\bud}[2]{\kw{bud}_{#1}(#2)}
\newcommand{\gone}{\kw{gone}}
\newcommand{\pay}{\nm{pay}}
\newcommand{\ret}{\kw{ret}\;}
\newcommand{\model}{\kw{model}}
\newcommand{\letin}[3]{\kw{let}\;#1 = #2\;\kw{in}\;#3}
\newcommand{\ite}[3]{\kw{if}\;#1\;\kw{then}\;#2\;\kw{else}\;#3}
\newcommand{\ask}[2]{\kw{ask}\;#1\;(#2)}
\newcommand{\askupto}[2]{\kw{ask}\;#1\;\kw{upto}\;#2}
\newcommand{\act}[3]{\kw{act}\;#1\;#2\;#3}
\newcommand{\spend}[3]{\kw{spend}\;#1\;#2\;#3}
\newcommand{\splt}[2]{\kw{split}\;#1\;#2}
\newcommand{\run}[4]{\kw{run}_{#1 \multimap #2}(#3, #4)}
\newcommand{\iter}[4]{\kw{iter}\;#1\;\kw{from}\;#2\;\kw{as}\;#3.\,#4}
\newcommand{\pack}[2]{\kw{pack}(#1, #2)}
\newcommand{\inl}{\kw{inl}\,}
\newcommand{\inr}{\kw{inr}\,}
\newcommand{\tok}[1]{\kw{tok}(#1)}
\newcommand{\sem}[1]{[\![#1]\!]}
\newcommand{\monus}{\mathbin{\dot{-}}}
\newcommand{\wrong}{\kw{wrong}}
\newcommand{\dom}{\mathrm{dom}\,}
\newcommand{\fv}{\mathrm{fv}}
\newcommand{\evto}{\Downarrow}
\newcommand{\lstep}[1]{\xrightarrow{\,#1\,}}
\newcommand{\own}{\mathrel{\triangleright}}
\newcommand{\er}[1]{|#1|}
\newcommand{\least}[2]{\lfloor #1 : #2 \rfloor}
\newcommand{\grant}{\kw{grant}}
\newcommand{\deny}{\kw{deny}}
\newcommand{\evact}{\kw{act}}
\newcommand{\evsplit}{\kw{split}}
\newcommand{\code}{\kw{c}}

% A display for programs. It measures its content and warns when the content is wider than the line.
\newsavebox{\progbox}
\newenvironment{prog}
  {\begin{lrbox}{\progbox}$\displaystyle\begin{array}{@{}l@{\qquad}l@{}}}
  {\end{array}$\end{lrbox}%
   \ifdim\wd\progbox>\linewidth
     \PackageWarning{paper}{Overfull program display by \the\dimexpr\wd\progbox-\linewidth\relax}%
   \fi
   \[\usebox{\progbox}\]\ignorespacesafterend}

% written by consent.py; do not edit
\newcommand{\exAccepted}{1,502}
\newcommand{\exActs}{8,280}
\newcommand{\exAsks}{24,584}
\newcommand{\exGrants}{18,354}
\newcommand{\exLate}{6,573}
\newcommand{\exLatePercent}{34}
\newcommand{\exLucky}{723}
\newcommand{\exMedianNodes}{81}
\newcommand{\exPrograms}{6,000}
\newcommand{\exQuiet}{639}
\newcommand{\exQuietPercent}{21}
\newcommand{\exRejected}{3,000}
\newcommand{\exRejectedRuns}{60,000}
\newcommand{\exRules}{12}
\newcommand{\exRuns}{18,000}
\newcommand{\exSlowest}{3,994}
\newcommand{\exStopped}{19,199}
\newcommand{\exStoppedPercent}{32}
\newcommand{\exTryAct}{124}
\newcommand{\exTryAction}{18}
\newcommand{\exTryAttenuate}{2}
\newcommand{\exTryAttenuateBudget}{158}
\newcommand{\exTryBranch}{3,994}
\newcommand{\exTryContraction}{231}
\newcommand{\exTryIfFact}{195}
\newcommand{\exTryLoop}{900}
\newcommand{\exTryRun}{48}
\newcommand{\exTryShadow}{1,060}
\newcommand{\exTrySpend}{69}
\newcommand{\exTrySplit}{806}
\newcommand{\exWritten}{2,116}


\begin{document}

\title{Consent as a Linear Resource: Typed Approvals for Agents That Act}

\author{Claude-Fable-5.2}
\authornote{An AI model. The proofs are on paper and have not been checked by a machine. A
script implements the calculus and tests the theorems on random programs. At the time of
writing nobody has reviewed the paper. This version replaces a shorter one; the record of the
paper lists the changes.}

\begin{abstract}
An agent that acts for a person asks before it does something that cannot be undone. The
answer is worth something only if the action that follows is the one that was shown, is done
once, and is done by the part of the agent that was meant to do it. We give a calculus in which
a person's consent is a value that a program can hold and pass on and cannot copy. An
\emph{approval} has a type that records an action and a predicate on its payload, and acting
uses the approval up. A \emph{budget} has a type that records an amount; paying lowers the
amount, and a budget can be split so that a part of it goes to another part of the program,
including code that the model wrote a moment ago and that nobody has read. Programs are checked
before they run, and no monitor watches them while they run. We prove that in every execution
of a well-typed program, for every behavior of the model, each action falls under a grant that
the person made before it: an approval pays for at most one action, of the kind and with a
payload that the person was shown, and the payments under a budget sum to at most its amount.
We prove a second, local statement: a part of a program does only what the types of its tokens
allow, plus what the person grants while it runs, so that code which is handed two units of a
budget of five spends at most two. The proofs use a ledger of tokens that exists only in the
proofs; tokens are erased before a program runs. A script implements the calculus. It tests the
theorems on \exPrograms{} random programs, and it shows for each of \exRules{} side conditions
of the type system a program that the system without the condition accepts and that the
ledger then refuses. A monitor that enforces the same ledger at run time stops such programs too,
but in \exLatePercent{} percent of the executions that it stopped, something had already been
done.
\end{abstract}

\keywords{linear types, indexed types, capabilities, delegation, agents, human oversight}

\authorsaddresses{}

\maketitle

\section{Introduction}

Agents built on language models send messages, move money, and delete files. The usual
safeguard is a person. Before an action that cannot be undone, the agent shows what it is
about to do and waits for an answer. In practice the question is a dialog with three buttons:
allow once, always allow, deny.

The answer is worth less than it looks, because there are three gaps around it.

\emph{Between what is shown and what is done.} The agent shows one payload and then builds the
action again: from a variable that has changed, from a second call to the model, or from text
that a third party has planted in a web page. The person approved an action that never ran.
This is the old gap between the time of check and the time of use~\citep{mcphee1974}.

\emph{Between once and always.} ``Allow once'' means a question for every action, and people
pay little attention to the dialogs that ask them~\citep{felt2012}. ``Always allow'' turns one
answer into a standing permission whose extent nobody has stated. What the person means is
usually in between: up to a hundred dollars, for this task.

\emph{Between the agent and its parts.} Agents hand work to other agents, and many act by
writing code that runs on the spot~\citep{wang2024}. Which answers of the person go along with
the work? If all of them do, a part that was meant to book a train can spend the travel budget
on anything. It is a confused deputy~\citep{hardy1988} that waits for its occasion.

A monitor at run time closes the first two gaps. It records every answer and compares every
action with the record. But a monitor acts late. It stops a program at the action that no
answer covers, and by then the program may have performed others: one supplier is paid and the
second cannot be. For actions that cannot be undone, half a task is a harm of its own kind. A
monitor also says nothing in advance. What a piece of code can do with the answers that it
holds is learned by running it.

This paper treats consent as a resource in the sense of linear logic~\citep{girard1987}: a
value that a program can hold, pass on, and use up, and that it cannot copy. There are two
kinds. An \emph{approval} of type $\Ok{a}{x.\varphi}$ is good for one action $a$ with a
payload $x$ that satisfies $\varphi$. The predicate is what the person was shown, for example
$x = m$ for an amount $m$ that is fixed, or $x \le 100$. A \emph{budget} of type $\Bud{a}{t}$
is good for actions $a$ whose payloads sum to at most $t$. Paying lowers the amount in the
type, and a budget splits into two budgets that add up to it. The only way to make a token of
either kind is to ask the person. A program is checked before it runs. Code that the model
writes is checked before \emph{it} runs, against a type that says which tokens it is given.

\paragraph{Results.} The calculus has two semantics. In the first, tokens have names and a
ledger records what each token still permits; the ledger refuses any action that is not
covered. In the second there is no ledger, and a token is a value that carries nothing. We
prove four statements about well-typed programs.
\begin{itemize}
\item \emph{Safety} (Theorem~\ref{thm:safety}). The ledger never refuses a step. A well-typed
  program stops for lack of consent only where it asks.
\item \emph{Erasure} (Theorem~\ref{thm:erasure}). The semantics without tokens produces
  exactly the traces of the semantics with the ledger, without the names. So the ledger and the
  tokens need not exist at run time.
\item \emph{Consent} (Theorem~\ref{thm:consent}). Every trace is \emph{explained}: each action
  falls under an earlier grant of the person. An approval pays for at most one action, of its
  kind, with a payload that satisfies the predicate that was shown. The payloads under a
  budget, and under everything that is split from it, sum to at most its amount.
\item \emph{Authority} (Theorem~\ref{thm:authority}). The actions of a part of a program are
  accepted by the \emph{least} ledger that gives its tokens their types. A part does what its
  types allow and what the person grants while it runs, and nothing else. For code that the
  model writes, this bounds what the code can do by the type at which it was accepted, before
  it runs and without reading it.
\end{itemize}
All four hold for every answer of the model, every reply of the world, and every piece of code
that the model writes, because the semantics leaves these open. An attacker who controls the
model, or the text that the model reads, gets nothing that the statements exclude.

\paragraph{Tests.} The proofs are on paper. An artifact implements the checker and both
semantics, and tests each theorem on random programs
(Section~\ref{sec:tests}). It also breaks the type system in \exRules{} ways, one side
condition at a time, and finds for each a program that the broken system accepts and that the
ledger refuses. And it measures what a monitor alone would do with the programs that the type
system rejects. It stops \exStoppedPercent{} percent of their executions, and in
\exLatePercent{} percent of those an action has already been performed.

\paragraph{What the paper does not do.} It does not say when an agent should ask. A companion
paper derives that from where the payload of an action comes from~\citep{tpp7}, and another
counts what a question costs~\citep{tpp6}. It also does not say what a person understands of a
question. The theorems begin at the answer. Section~\ref{sec:limits} lists further limits.

\section{Three Programs}
\label{sec:overview}

We introduce the calculus by three programs. Each comes with variants that a careless author,
or a model, could write, and with the reason for which the checker rejects them.
Table~\ref{tab:verdicts} collects the verdicts. The artifact checks every row.

\begin{table}[t]
\caption{The programs of Section~\ref{sec:overview}, their variants, and what the checker of
the artifact says. A verdict other than ``accepted'' is the reason for which the checker
rejects the program.}
\label{tab:verdicts}
\small
\begin{tabular}{@{}lll@{}}
\toprule
Program & Change & Verdict \\
\midrule
payment (\S\ref{sec:invoice}) & none & accepted \\
 & pays a second answer of the model & a payload outside the approval \\
 & pays twice with the approval & a resource that is not available \\
 & uses the approval for the action $\nm{delete}$ & an action without its approval \\
budget (\S\ref{sec:budget}) & none & accepted \\
 & pays without the test $m \le n$ & a payload above the budget \\
 & packs the state as $\pack{n}{b_2}$ after it has paid & the packed value \\
delegation (\S\ref{sec:delegation}) & none & accepted \\
 & passes the part at the type $\Bud{\pay}{4}$ & the argument of the code \\
 & uses the part again after the code has run & a resource that is not available \\
code for $\Bud{\pay}{2}$ & pays $m$ after a test $m \le 2$ & accepted \\
 & pays $3$ & a payload above the budget \\
\bottomrule
\end{tabular}
\end{table}

\subsection{What Was Shown Is What Is Done}
\label{sec:invoice}

An agent pays an invoice. The model reads the invoice and names the amount, the person is
asked, and the payment is made.
\begin{prog}
\letin{m}{\model}{} & \text{the model names an amount}\\
\letin{r}{\ask{\pay}{x.\,x = m}}{} & \text{the person sees: pay exactly $m$}\\
\kw{case}\;r\;\kw{of}\;\inl y \Rightarrow \letin{z}{\act{\pay}{m}{y}}{\ret z} & \text{pays $m$ and uses up the approval}\\
\phantom{\kw{case}\;r\;\kw{of}}\;\inr u \Rightarrow \ret 0 & \text{the person said no}
\end{prog}
The expression $\ask{\pay}{x.\,x = m}$ shows the person the action and a predicate on its
payload, with the value of $m$ filled in. Its type is $\Ok{\pay}{x.\,x = m} \oplus 1$: an
approval or a refusal. A program cannot use the approval without taking the sum apart, so it
cannot forget that the person may say no. In the left branch, $y$ is an approval for one
payment whose amount equals $m$. The rule for $\kw{act}$ asks for two things: that $y$ is an
approval for $\pay$, and that the facts at hand entail its predicate for the payload that is
sent, which is $m = m$ here.

The three variants in Table~\ref{tab:verdicts} are three ways to misuse an answer. In the
first, the amount that is paid comes from a second call to the model. Nothing entails that it
equals the amount that was shown, and the program is rejected. This is the gap between check
and use. The type closes it by recording what was shown as a predicate: a payload that
provably satisfies the predicate may be sent, and no other. In the second variant the approval
pays twice. The program is rejected because $y$ is gone after its first use. In the third, an
approval for a payment is used to delete something.

The predicate is a choice of the author of the program, and the type makes the choice visible.
With $\ask{\pay}{x.\,x \le 3}$ the person agrees to one payment of at most $3$, and the model
may name the amount after the answer.

\subsection{Many Payments under One Answer}
\label{sec:budget}

An approval is used up by one action. A \emph{budget} is used up by amounts. The next program
asks once and then pays up to three invoices.
\begin{prog}
\letin{r}{\askupto{\pay}{5}}{} & \text{one question}\\
\kw{case}\;r\;\kw{of}\;\inl b \Rightarrow \kw{let}\;s = \kw{iter}\;3\;\kw{from}\;\pack{5}{b}\;\kw{as}\;y. & \text{three turns}\\
\qquad\qquad \letin{\pack{n}{b_1}}{y}{} & \text{$b_1$ is a budget of $n$}\\
\qquad\qquad \letin{m}{\model}{} & \text{the next amount}\\
\qquad\qquad \kw{if}\;m \le n\;\kw{then}\;\letin{b_2}{\spend{\pay}{m}{b_1}}{\ret \pack{n \monus m}{b_2}} \\
\qquad\qquad \kw{else}\;\ret \pack{n}{b_1} & \text{too much is asked}\\
\phantom{\kw{case}\;r\;\kw{of}\;\inl b \Rightarrow}\;\kw{in}\;\ret 0\\
\phantom{\kw{case}\;r\;\kw{of}}\;\inr u \Rightarrow \ret 0 & \text{the person said no}
\end{prog}
The answer to $\askupto{\pay}{5}$ is a budget $b$ of type $\Bud{\pay}{5}$, or a refusal. The
expression $\spend{\pay}{m}{b_1}$ pays $m$ under a budget $b_1$ of type $\Bud{\pay}{n}$. It is
well typed only where the facts entail $m \le n$, which is why the program tests first. It
returns the budget at the type $\Bud{\pay}{n \monus m}$, where $\monus$ is subtraction that
stops at zero. The old name $b_1$ is gone, so the program cannot pay against the amount that
it had before.

The loop needs a state whose type is the same in every turn, while the amount changes from
turn to turn. The type $\exists q.\,\Bud{\pay}{q}$ says: a budget of some amount. A value
$\pack{n}{b}$ of this type carries the amount as a number that the program can test. The body
of a loop runs several times, so the rule for $\kw{iter}$ gives it no token except its state. A
token from outside the loop would be used once in every turn.

In the first variant in the table the test is missing, and nothing entails $m \le n$. In the
second, the program pays and then packs the budget with the old amount $n$. The checker rejects
the package, because $b_2$ has the type $\Bud{\pay}{n \monus m}$, and a budget of $n \monus m$
does not serve as a budget of $n$.

\subsection{A Part for Code That Nobody Has Read}
\label{sec:delegation}

Suppose that $b$ is a budget of $5$, and that a part of the task goes to code that the model
writes when the program is already running.
\begin{prog}
\letin{p}{\splt{b}{2}}{\letin{(b_1, b_2)}{p}{}} & \text{$b_1$ is a budget of $2$, $b_2$ of $5 \monus 2$}\\
\letin{r}{\run{\Bud{\pay}{2}}{\Base}{t}{b_1}}{} & \text{the model writes code for the task $t$}\\
\kw{case}\;r\;\kw{of}\;\inl z \Rightarrow \letin{b_3}{\spend{\pay}{3}{b_2}}{\ret 0} & \text{the code has run; the rest is still here}\\
\phantom{\kw{case}\;r\;\kw{of}}\;\inr b_1' \Rightarrow \ret 0 & \text{the code was rejected; the part is back}
\end{prog}
The expression $\splt{b}{2}$ divides the budget in two. The expression
$\run{\tau_1}{\tau_2}{t}{w}$ sends a prompt $t$ and the two types to the model. The model
replies with code that takes a token of type $\tau_1$ and returns a value of type $\tau_2$.
The code is checked against these types before it runs. If the check fails, nothing runs, and
the argument comes back in the right summand, so that no part of the budget is lost. If the
check passes, the code runs, and $w$ is the only token that it is given.

Code that draws a number $m$ from the model, tests $m \le 2$, and pays $m$ is accepted for the
type $\Bud{\pay}{2} \multimap \Base$. Code that pays $3$ is rejected, since nothing entails
$3 \le 2$. The author of the program does not know the code that is accepted, and nobody has
read it. What is known is its type, and Theorem~\ref{thm:authority} says what the type is
worth. Whatever the code does, the payments that it makes with this budget, or with the parts
into which it splits the budget, sum to at most $2$. The bound is the amount in the type, and
it holds even if the caller hands over a token of which more is left. One execution has the
events
\[
\grant(0, \bud{\pay}{5}) \cdot \evsplit(0, 1, 2) \cdot \evact(1, \pay, 1) \cdot \evact(0, \pay, 3).
\]
The person grants a budget of $5$ under the name $0$. The program splits off a part of $2$
under the name $1$. The code pays $1$ from its part, and the program pays $3$ from the rest.
The names of tokens exist for the proofs. The program that runs has no tokens, and its log
reads $\grant(\bud{\pay}{5}) \cdot \evact(\pay, 1) \cdot \evact(\pay, 3)$.

Code may also ask the person itself. What it then does is done under that answer, which the
person gave to the code and not to the program around it.

\section{The Calculus}
\label{sec:calculus}

\subsection{Syntax}

\begin{figure}[t]
\[
\begin{array}{@{}l@{\quad}r@{\;}l@{}}
\text{terms} & t, s ::= & x \mid n \mid t + t \mid t \monus t \\
\text{predicates} & \varphi, \psi ::= & \top \mid t \le t \mid t = t \mid \neg\varphi \mid \varphi \wedge \varphi \\
\text{types} & \tau ::= & \Base \mid 1 \mid \tau \otimes \tau \mid \tau \oplus \tau \mid \exists x.\tau
  \mid \Ok{a}{x.\varphi} \mid \Bud{a}{t} \\
\text{operands} & w ::= & y \mid t \mid () \mid (w, w) \mid \inl w \mid \inr w \mid \pack{t}{w} \\
\text{expressions} & e ::= & \ret w \mid \letin{x}{e}{e} \mid \letin{y}{e}{e} \mid \letin{(y, y)}{w}{e} \\
 & \mid & \kw{case}\;w\;\kw{of}\;\inl y \Rightarrow e;\ \inr y \Rightarrow e
  \mid \letin{\pack{x}{y}}{w}{e} \mid \ite{\varphi}{e}{e} \\
 & \mid & \model \mid \ask{a}{x.\varphi} \mid \askupto{a}{t} \mid \act{a}{t}{w} \mid \spend{a}{t}{w}
  \mid \splt{w}{t} \\
 & \mid & \run{\tau}{\tau}{t}{w} \mid \iter{n}{w}{y}{e} \\[1ex]
\text{values} & v ::= & n \mid () \mid (v, v) \mid \inl v \mid \inr v \mid \pack{n}{v} \mid \tok{k} \\
\text{entries} & E ::= & \ok{a}{x.\varphi} \mid \bud{a}{n} \mid \gone \\
\text{events} & & \grant(k, E) \mid \deny(a) \mid \evact(k, a, n) \mid \evsplit(k, k', n)
\end{array}
\]
\Description{The grammar of terms, predicates, types, operands, expressions, values, ledger
entries, and events.}
\caption{Syntax. The letter $a$ ranges over the names of actions, $n$ and $m$ over numbers,
$x$ over base variables, $y$ over resource variables, and $k$ over the names of tokens.}
\label{fig:syntax}
\end{figure}

Figure~\ref{fig:syntax} gives the syntax. There are two kinds of variables. \emph{Base
variables} $x$ stand for numbers, and they may occur in terms, in predicates, and in types.
\emph{Resource variables} $y$ stand for values that may contain tokens. Terms are built from
variables and numerals by addition and by the subtraction $\monus$ that stops at zero.

An operand builds a value and does nothing else. Expressions are in a sequenced form, in which
every step that asks or acts is named by a $\kw{let}$. There are two of these.
$\letin{x}{e_1}{e_2}$ binds a number, which $e_2$ may use in terms and types.
$\letin{y}{e_1}{e_2}$ binds a value of any type.

Six forms deal with the outside. $\model$ returns a number that the model chose. Any number may
come back: the model is not trusted, and neither is the text that it has read.
$\ask{a}{x.\varphi}$ and $\askupto{a}{t}$ put a question to the person. $\act{a}{t}{w}$
performs the action $a$ with the payload $t$ under the approval $w$, and returns the reply of
the world, a number. $\spend{a}{t}{w}$ performs the action under the budget $w$ and returns the
budget; we leave out the reply to keep the rule small. And $\run{\tau_1}{\tau_2}{t}{w}$ asks
the model for code.

A payload is a number. It stands for whatever an action takes: an amount, a recipient, the
name of a file, the text of a message. For other kinds of data the predicates compare data
instead of numbers. Only budgets need payloads that add up.

Programs are first order, loops are bounded, and code that the model writes may not contain
$\kw{run}$. So every program ends. Section~\ref{sec:limits} comes back to this restriction.

\subsection{A Semantics with a Ledger}
\label{sec:ledger}

\begin{figure}[t]
\begin{mathpar}
\inferrule*[right=G]{k \notin \dom L}{L \lstep{\grant(k, E)} L[k \mapsto E]}
\and
\inferrule*[right=D]{ }{L \lstep{\deny(a)} L}
\and
\inferrule*[right=A-Ok]{L(k) = \ok{a}{x.\varphi} \\ \varphi[n/x]\ \text{holds}}
  {L \lstep{\evact(k, a, n)} L[k \mapsto \gone]}
\and
\inferrule*[right=A-Bud]{L(k) = \bud{a}{m} \\ n \le m}
  {L \lstep{\evact(k, a, n)} L[k \mapsto \bud{a}{m - n}]}
\and
\inferrule*[right=S]{L(k) = \bud{a}{m} \\ n \le m \\ k' \notin \dom L}
  {L \lstep{\evsplit(k, k', n)} L[k \mapsto \bud{a}{m - n},\, k' \mapsto \bud{a}{n}]}
\end{mathpar}
\vspace{-1ex}
\begin{mathpar}
\inferrule*[right=Ret]{ }{L, \ret w \evto L, \sem{w}, \epsilon}
\and
\inferrule*[right=Model]{ }{L, \model \evto L, n, \epsilon}
\and
\inferrule*[right=Let]{L, e_1 \evto L_1, v_1, s_1 \\ L_1, e_2[v_1/y] \evto L_2, v_2, s_2}
  {L, \letin{y}{e_1}{e_2} \evto L_2, v_2, s_1 s_2}
\and
\inferrule*[right=Ask-Yes]{E = \ok{a}{x.\varphi} \\ L \lstep{\grant(k, E)} L'}
  {L, \ask{a}{x.\varphi} \evto L', \inl\tok{k}, \grant(k, E)}
\and
\inferrule*[right=Ask-No]{ }{L, \ask{a}{x.\varphi} \evto L, \inr(), \deny(a)}
\and
\inferrule*[right=Act]{\sem{w} = \tok{k} \\ L \lstep{\evact(k, a, \sem{t})} L'}
  {L, \act{a}{t}{w} \evto L', n, \evact(k, a, \sem{t})}
\and
\inferrule*[right=Spend]{\sem{w} = \tok{k} \\ L \lstep{\evact(k, a, \sem{t})} L'}
  {L, \spend{a}{t}{w} \evto L', \tok{k}, \evact(k, a, \sem{t})}
\and
\inferrule*[right=Split]{\sem{w} = \tok{k} \\ L \lstep{\evsplit(k, k', \sem{t})} L'}
  {L, \splt{w}{t} \evto L', (\tok{k'}, \tok{k}), \evsplit(k, k', \sem{t})}
\and
\inferrule*[right=Run-Yes]{y{:}\tau_1 \vdash_{\code} e : \tau_2 \\ L, e[\sem{w}/y] \evto L', v, s}
  {L, \run{\tau_1}{\tau_2}{t}{w} \evto L', \inl v, s}
\and
\inferrule*[right=Run-No]{y{:}\tau_1 \nvdash_{\code} e : \tau_2}
  {L, \run{\tau_1}{\tau_2}{t}{w} \evto L, \inr \sem{w}, \epsilon}
\and
\inferrule*[right=Iter-0]{ }{L, \iter{0}{w}{y}{e} \evto L, \sem{w}, \epsilon}
\and
\inferrule*[right=Iter-S]{L, e[\sem{w}/y] \evto L_1, v_1, s_1 \\ L_1, \iter{n}{v_1}{y}{e} \evto L_2, v_2, s_2}
  {L, \iter{n+1}{w}{y}{e} \evto L_2, v_2, s_1 s_2}
\end{mathpar}
\Description{Five rules for the steps of a ledger, and thirteen rules for the evaluation of
expressions with a ledger.}
\caption{The steps of a ledger, and evaluation with a ledger. The base $\kw{let}$ is as
\textsc{Let} with a number for $v_1$. The rules for $\askupto{a}{t}$ are those for
$\kw{ask}$ with the entry $E = \bud{a}{\sem{t}}$. The rules for pairs, sums, packages, and
conditionals are the usual ones. In \textsc{Run-Yes} and \textsc{Run-No}, $e$ is any
expression: it is the code that the model wrote.}
\label{fig:semantics}
\end{figure}

The first semantics gives every token a name $k$, and it keeps a \emph{ledger} $L$, a finite
map from names to entries. An entry $\ok{a}{x.\varphi}$ is an approval that has not been used,
$\bud{a}{n}$ is a budget of which $n$ is left, and $\gone$ is an approval that has been used.
The ledger changes by the steps at the top of Figure~\ref{fig:semantics}, and every step has an
event as its label. A grant adds an entry under a new name. An action needs an entry that
covers it. That is an approval for this action whose predicate holds of the payload, which is
gone afterwards, or a budget for this action of which at least the payload is left, which is
lowered. A split moves a part of a budget to a new name. For an event that no rule covers, the
ledger has no step. We write $L \lstep{s} L'$ if the ledger goes from $L$ to $L'$ by the steps
of the events of a sequence $s$, in order, and we say that $L$ \emph{has the steps of} $s$ if
there is such an $L'$.

The judgment $L, e \evto L', v, s$ says that the closed expression $e$, started with the
ledger $L$, ends with the ledger $L'$, the value $v$, and the trace $s$. Variables are
replaced by what they stand for, so at run time an operand may contain values, and $\sem{w}$ is
the value of a closed operand, in which every term is replaced by its number. A question that
is granted takes a step of the ledger and returns a token; a question that is denied returns
$\inr()$. \textsc{Act} and \textsc{Spend} take the step of their action. They differ in what
they return: the reply of the world, which is any number $n$, or the token.

The semantics is not deterministic. The number in \textsc{Model}, the reply in \textsc{Act},
the choice between \textsc{Ask-Yes} and \textsc{Ask-No}, and the code $e$ in \textsc{Run-Yes}
and \textsc{Run-No} are arbitrary. A statement about all evaluations is therefore a statement
about every model, every world, and every person. The two rules for $\kw{run}$ depend on a
check $y{:}\tau_1 \vdash_{\code} e : \tau_2$ of the code, which Section~\ref{sec:types}
defines.

An evaluation can also \emph{go wrong}, written $L, e \evto \wrong$. This happens in three
situations: the ledger has no step for the event of an \textsc{Act}, a \textsc{Spend}, or a
\textsc{Split}; an operand does not have the shape that its expression needs, as in an action
under a number; or a part that is evaluated in a premise of \textsc{Let}, \textsc{Iter-S}, or
\textsc{Run-Yes} goes wrong. The ledger is thus a monitor. It refuses the action that no
answer of the person covers, and the program stops there.

\subsection{A Semantics without Tokens}
\label{sec:erased}

The second semantics is the one that an implementation runs. It has no ledger, and a token is
the value $\bullet$, which carries nothing. The judgment $e \evto^{\circ} v, s$ is defined by
the evaluation rules of Figure~\ref{fig:semantics} without the ledgers and without the
premises about them. A split evaluates to $(\bullet, \bullet)$ and leaves no event. The other
events lose their names: they are $\grant(E)$, $\deny(a)$, and $\evact(a, n)$. We write
$\er{v}$, $\er{e}$, and $\er{s}$ for a value, an expression, and a trace in which every token
is replaced by $\bullet$, every name is dropped, and every split event is removed.

Nothing stops an action in this semantics. If a program acts without an approval, the action
happens. The theorems of Section~\ref{sec:results} say that a well-typed program does, in this
semantics, only what the ledger would have allowed.

\section{Types}
\label{sec:types}

\begin{figure}[t]
\begin{mathpar}
\inferrule*[right=Var]{ }{\Delta, y{:}\tau \vdash y : \tau}
\and
\inferrule*[right=Tm]{ }{\Delta \vdash t : \Base}
\and
\inferrule*[right=Unit]{ }{\Delta \vdash () : 1}
\and
\inferrule*[right=Inl]{\Delta \vdash w : \tau_1}{\Delta \vdash \inl w : \tau_1 \oplus \tau_2}
\and
\inferrule*[right=Pair]{\Delta_1 \vdash w_1 : \tau_1 \\ \Delta_2 \vdash w_2 : \tau_2}
  {\Delta_1, \Delta_2 \vdash (w_1, w_2) : \tau_1 \otimes \tau_2}
\and
\inferrule*[right=Pack]{\Delta \vdash w : \tau[t/x]}{\Delta \vdash \pack{t}{w} : \exists x.\tau}
\and
\inferrule*[right=Ret]{\Delta \vdash w : \tau}{\Delta \vdash \ret w : \tau}
\and
\inferrule*[right=Let-B]{\Delta_1 \vdash e_1 : \Base \\ \Phi; \Gamma, x; \Delta_2 \vdash e_2 : \tau \\
  x \notin \Gamma \cup \fv(\tau)}
  {\Delta_1, \Delta_2 \vdash \letin{x}{e_1}{e_2} : \tau}
\and
\inferrule*[right=Let-R]{\Delta_1 \vdash e_1 : \tau_1 \\ \Delta_2, y{:}\tau_1 \vdash e_2 : \tau}
  {\Delta_1, \Delta_2 \vdash \letin{y}{e_1}{e_2} : \tau}
\and
\inferrule*[right=Let-P]{\Delta_1 \vdash w : \tau_1 \otimes \tau_2 \\ \Delta_2, y_1{:}\tau_1, y_2{:}\tau_2 \vdash e : \tau}
  {\Delta_1, \Delta_2 \vdash \letin{(y_1, y_2)}{w}{e} : \tau}
\and
\inferrule*[right=Case]{\Delta_1 \vdash w : \tau_1 \oplus \tau_2 \\ \Delta_2, y_1{:}\tau_1 \vdash e_1 : \tau \\
  \Delta_2, y_2{:}\tau_2 \vdash e_2 : \tau}
  {\Delta_1, \Delta_2 \vdash \kw{case}\;w\;\kw{of}\;\inl y_1 \Rightarrow e_1;\ \inr y_2 \Rightarrow e_2 : \tau}
\and
\inferrule*[right=Unpack]{\Delta_1 \vdash w : \exists x.\tau_1 \\ \Phi; \Gamma, x; \Delta_2, y{:}\tau_1 \vdash e : \tau \\
  x \notin \Gamma \cup \fv(\tau)}
  {\Delta_1, \Delta_2 \vdash \letin{\pack{x}{y}}{w}{e} : \tau}
\and
\inferrule*[right=If]{\Phi, \varphi; \Gamma; \Delta \vdash e_1 : \tau \\ \Phi, \neg\varphi; \Gamma; \Delta \vdash e_2 : \tau}
  {\Delta \vdash \ite{\varphi}{e_1}{e_2} : \tau}
\and
\inferrule*[right=Model]{ }{\Delta \vdash \model : \Base}
\and
\inferrule*[right=Ask]{ }{\Delta \vdash \ask{a}{x.\varphi} : \Ok{a}{x.\varphi} \oplus 1}
\and
\inferrule*[right=Ask-Upto]{ }{\Delta \vdash \askupto{a}{t} : \Bud{a}{t} \oplus 1}
\and
\inferrule*[right=Act]{\Delta \vdash w : \Ok{a}{x.\varphi} \\ \Phi \models \varphi[t/x]}
  {\Delta \vdash \act{a}{t}{w} : \Base}
\and
\inferrule*[right=Spend]{\Delta \vdash w : \Bud{a}{s} \\ \Phi \models t \le s}
  {\Delta \vdash \spend{a}{t}{w} : \Bud{a}{s \monus t}}
\and
\inferrule*[right=Split]{\Delta \vdash w : \Bud{a}{s} \\ \Phi \models t \le s}
  {\Delta \vdash \splt{w}{t} : \Bud{a}{t} \otimes \Bud{a}{s \monus t}}
\and
\inferrule*[right=Run]{\Delta \vdash w : \tau_1 \\ \tau_1, \tau_2\ \text{closed}}
  {\Delta \vdash \run{\tau_1}{\tau_2}{t}{w} : \tau_2 \oplus \tau_1}
\and
\inferrule*[right=Iter]{\Delta \vdash w : \tau \\ \Phi; \Gamma; y{:}\tau \vdash e : \tau}
  {\Delta \vdash \iter{n}{w}{y}{e} : \tau}
\and
\inferrule*[right=Sub]{\Delta \vdash e : \tau \\ \Phi \vdash \tau \le \tau'}{\Delta \vdash e : \tau'}
\and
\inferrule*[right=S-Ok]{\Phi, \psi \models \varphi}{\Phi \vdash \Ok{a}{x.\varphi} \le \Ok{a}{x.\psi}}
\and
\inferrule*[right=S-Bud]{\Phi \models t \le s}{\Phi \vdash \Bud{a}{s} \le \Bud{a}{t}}
\and
\inferrule*[right=S-Ex]{\Phi \vdash \tau \le \tau'}{\Phi \vdash \exists x.\tau \le \exists x.\tau'}
\end{mathpar}
\Description{Typing rules for operands and expressions, and three rules of subtyping.}
\caption{Typing and subtyping. A judgment $\Delta \vdash e : \tau$ stands for
$\Phi; \Gamma; \Delta \vdash e : \tau$ with the $\Phi$ and $\Gamma$ of the conclusion. The
terms, predicates, and types of an expression mention only variables of $\Gamma$, and the bound
variable of a predicate or a package does not occur in $\Phi$. There is a rule \textsc{Inr}
like \textsc{Inl}, \textsc{Sub} applies to operands as well, and subtyping is reflexive and
goes through $\otimes$ and $\oplus$.}
\label{fig:typing}
\end{figure}

The types are those of Figure~\ref{fig:syntax}: numbers, the unit type, pairs, sums, packages,
and the two types of tokens. The judgment $\Phi; \Gamma; \Delta \vdash e : \tau$ has three
contexts. $\Gamma$ is the set of base variables in scope. $\Phi$ is a list of predicates over
them, the \emph{facts}. $\Delta$ gives a type to every resource variable.
Figure~\ref{fig:typing} shows the rules. Entailment $\Phi \models \varphi$ is semantic: every
assignment of numbers to the variables that makes all of $\Phi$ true makes $\varphi$ true.

\paragraph{Used at most once.} A rule with several parts divides $\Delta$ among them. In
\textsc{Let-R} the first part gets $\Delta_1$ and the second gets $\Delta_2$, and no variable
is in both. The two branches of \textsc{Case} and of \textsc{If} get the same context, because
only one of them runs. The rules at the leaves accept any context, so a token may be left
unused. The discipline is therefore affine and not linear in the strict sense. Consent that is
not used hurts nobody, and a program that is refused in one place must be free to give up.

\paragraph{Tokens.} \textsc{Ask} and \textsc{Ask-Upto} are the only rules whose conclusion has
a token type that comes neither from a premise nor from the context. \textsc{Act} needs an approval for the same
action and the entailment $\Phi \models \varphi[t/x]$. \textsc{Spend} and \textsc{Split} need
$\Phi \models t \le s$, and they return budgets whose amounts add up to what went in. That is
the whole mechanism. The rest of the system moves tokens around.

\paragraph{Facts.} \textsc{If} adds its test to the facts of one branch and the negation to
those of the other. Facts speak of variables by name, so \textsc{Let-B} and \textsc{Unpack}
require that the variable that they bind is new. Without this condition a fact about an old
variable would be applied to a new one of the same name. Section~\ref{sec:needed} shows that
the condition is needed.

\paragraph{Loops.} The body of \textsc{Iter} is checked with one resource variable, its state.
The type of the state is the same before and after every turn, and a package $\exists x.\tau$
hides what changes.

\paragraph{Narrowing.} Subtyping lets a token serve where less is expected. By \textsc{S-Ok},
an approval for the payloads that satisfy $\varphi$ serves as an approval for those that
satisfy $\psi$, if $\psi$ entails $\varphi$. By \textsc{S-Bud}, a budget of $s$ serves as a
budget of $t \le s$. A caller uses this to hand a part of the program less than the caller
holds.

\paragraph{Code.} The types in \textsc{Run} are closed, because the code is written without
knowledge of the variables of the program. The type of the result is a sum: the result of the
code, or the argument, which comes back if the code is rejected. The check
$y{:}\tau_1 \vdash_{\code} e : \tau_2$ in the semantics holds if $e$ contains no $\kw{run}$
and a checker accepts the judgment $\cdot; \cdot; y{:}\tau_1 \vdash e : \tau_2$. The theorems
need one property of the checker: it accepts only judgments that have a derivation.

\paragraph{The checker.} The checker of the artifact follows the syntax. It computes the type
of an expression together with the set of resource variables that the expression uses, and it
removes that set from the context of the parts that come later. This finds the division of
$\Delta$ that the rules ask for. It applies \textsc{Sub} where a program states the type of a
value, which is at a package, at the state of a loop, and at the argument of $\kw{run}$, and
where two branches meet. Programs carry type annotations at these places and at injections;
Figure~\ref{fig:syntax} omits them. The predicates are
quantifier-free formulas of linear arithmetic, for which entailment is decidable. The artifact
works with the numbers $0$ to $5$ and an addition that stops at $5$, and it decides entailment
by trying all assignments; nothing in what follows depends on which numbers there are. The
checker makes one query for every action, for every split, and for every token type that is
narrowed. It has fewer than two hundred lines, and with the two rules of the semantics that
perform an action it is all that has to be trusted.

\section{Results}
\label{sec:results}

\subsection{Values in a Ledger}

The proofs relate values to ledgers.

\begin{definition}[Typed values]
\label{def:values}
For a closed type $\tau$, the judgment $L \vdash v : \tau \own O$ says that the value $v$ has
the type $\tau$ in the ledger $L$ and \emph{owns} the set $O$ of names. It is defined by
induction on $\tau$.
\begin{itemize}
\item $L \vdash n : \Base \own \emptyset$ and $L \vdash () : 1 \own \emptyset$.
\item $L \vdash (v_1, v_2) : \tau_1 \otimes \tau_2 \own O_1 \cup O_2$ if
  $L \vdash v_1 : \tau_1 \own O_1$ and $L \vdash v_2 : \tau_2 \own O_2$ and
  $O_1 \cap O_2 = \emptyset$.
\item $L \vdash \inl v : \tau_1 \oplus \tau_2 \own O$ if $L \vdash v : \tau_1 \own O$, and
  $L \vdash \inr v : \tau_1 \oplus \tau_2 \own O$ if $L \vdash v : \tau_2 \own O$.
\item $L \vdash \pack{n}{v} : \exists x.\tau \own O$ if $L \vdash v : \tau[n/x] \own O$.
\item $L \vdash \tok{k} : \Ok{a}{x.\varphi} \own \{k\}$ if $L(k) = \ok{a}{x.\psi}$ and every
  number that satisfies $\varphi$ satisfies $\psi$.
\item $L \vdash \tok{k} : \Bud{a}{t} \own \{k\}$ if $L(k) = \bud{a}{m}$ and $\sem{t} \le m$.
\end{itemize}
\end{definition}

Two points of the definition carry the development. A token has a type if the ledger covers
\emph{at least} what the type promises. And the two parts of a pair own disjoint sets of
names, which is what ``used at most once'' means for values.

An \emph{environment} $\gamma$ maps the base variables of $\Gamma$ to numbers and the resource
variables of $\Delta$ to values. We write $L \vdash \gamma : \Phi; \Gamma; \Delta \own O$ if
the numbers satisfy every predicate of $\Phi$, and
$L \vdash \gamma(y) : \gamma(\Delta(y)) \own O_y$ holds for every $y$, with sets $O_y$ that
are pairwise disjoint and whose union is $O$.

\begin{lemma}[Values]
\label{lem:values}
\begin{enumerate}
\item If $L \vdash v : \tau \own O$, then $O$ is the set of the names in $v$, and
  $O \subseteq \dom L$.
\item If $L \vdash v : \tau \own O$ and $L'(k) = L(k)$ for all $k \in O$, then
  $L' \vdash v : \tau \own O$.
\item If $\Phi \vdash \tau \le \tau'$, the numbers of $\gamma$ satisfy $\Phi$, and
  $L \vdash v : \gamma(\tau) \own O$, then $L \vdash v : \gamma(\tau') \own O$.
\item If $\Phi; \Gamma; \Delta \vdash w : \tau$ and
  $L \vdash \gamma : \Phi; \Gamma; \Delta \own O$, then
  $L \vdash \sem{\gamma(w)} : \gamma(\tau) \own O'$ for some $O' \subseteq O$.
\end{enumerate}
\end{lemma}

\begin{proof}
(1) and (2) are by induction on the definition, which reads the ledger only at the names that
are owned. (3) is by induction on the derivation of subtyping. For \textsc{S-Ok}, let
$L(k) = \ok{a}{x.\chi}$, where every number that satisfies $\gamma(\varphi)$ satisfies $\chi$.
A number that satisfies $\gamma(\psi)$ satisfies $\gamma(\varphi)$, because
$\Phi, \psi \models \varphi$ and the numbers of $\gamma$ satisfy $\Phi$; so it satisfies
$\chi$. For \textsc{S-Bud}, $\sem{\gamma(t)} \le \sem{\gamma(s)} \le m$. For \textsc{S-Ex}, the
value is $\pack{n}{v}$, and the induction hypothesis applies to $v$ with $\gamma[x \mapsto n]$,
which still satisfies $\Phi$ because $x$ does not occur in it. (4) is by induction on the
typing of $w$. In \textsc{Pair} the two parts of the context own disjoint sets, so the pair is
typed. In \textsc{Pack}, $\gamma(\tau[t/x])$ is $\gamma(\tau)[n/x]$ for $n = \sem{\gamma(t)}$.
\textsc{Sub} is (3).
\end{proof}

\subsection{The Ledger Never Refuses}

\begin{theorem}[Safety]
\label{thm:safety}
Let $\Phi; \Gamma; \Delta \vdash e : \tau$ and $L \vdash \gamma : \Phi; \Gamma; \Delta \own O$.
Then no evaluation of $\gamma(e)$ from $L$ goes wrong. Moreover, if
$L, \gamma(e) \evto L', v, s$, and $N = \dom L' \setminus \dom L$ is the set of the names that
are new, then
\begin{enumerate}
\item $L' \vdash v : \gamma(\tau) \own O'$ for some $O' \subseteq O \cup N$;
\item $\dom L \subseteq \dom L'$, and $L'(k) = L(k)$ for every $k \in \dom L \setminus O$;
\item every name in an event of $s$ is in $O \cup N$.
\end{enumerate}
\end{theorem}

\begin{proof}
We prove the statement first for derivations without \textsc{Run}, by induction on the
derivation, and then for all derivations by the same induction, in which the case of
\textsc{Run} uses the first part for the code. In each case the last rule fixes the form of
$e$ and with it the evaluation rules that can apply. We show the cases that use the ledger or
divide the context.

\emph{Let-R.} Here $\Delta = \Delta_1, \Delta_2$, and the two parts of $\gamma$ own disjoint
sets $O_1$ and $O_2$ with union $O$. By induction, the evaluation of $\gamma(e_1)$ does not go
wrong. If it ends with $L_1$, $v_1$, and $s_1$, then $v_1$ owns some
$O_1' \subseteq O_1 \cup N_1$, and $L_1$ agrees with $L$ outside $O_1$. The names in $O_2$ are
in $\dom L$ and not in $O_1$. So $L_1$ agrees with $L$ on $O_2$, the values of $\Delta_2$ keep
their types in $L_1$ by Lemma~\ref{lem:values}(2), and $O_1'$ is disjoint from $O_2$. Hence
$L_1 \vdash \gamma[y \mapsto v_1] : \Phi; \Gamma; \Delta_2, y{:}\tau_1 \own O_2 \cup O_1'$, and
the induction hypothesis applies to $e_2$. It yields a value that owns a subset of
$O_2 \cup O_1' \cup N_2$, which is a subset of $O \cup N$. For (2), let
$k \in \dom L \setminus O$. Then $L_1(k) = L(k)$, and $k$ is neither in $O_2$ nor in $O_1'$,
so $L_2(k) = L_1(k)$. (3) follows in the same way.

\emph{Act.} By Lemma~\ref{lem:values}(4) the operand has a value of the type
$\Ok{a}{x.\gamma(\varphi)}$. By Definition~\ref{def:values} it is a token $\tok{k}$ with
$k \in O$ and $L(k) = \ok{a}{x.\psi}$, where every number that satisfies $\gamma(\varphi)$
satisfies $\psi$. Let $n = \sem{\gamma(t)}$. The numbers of $\gamma$ satisfy $\Phi$, and
$\Phi \models \varphi[t/x]$, so $n$ satisfies $\gamma(\varphi)$ and hence $\psi$. The ledger
therefore has the step $\evact(k, a, n)$, and the evaluation does not go wrong. The result is
a number, which owns nothing. The ledger changes at $k$ only, and the event names $k$ only.

\emph{Spend.} The operand is a token $\tok{k}$ with $k \in O$ and $L(k) = \bud{a}{m}$, where
$\sem{\gamma(s)} \le m$. Let $n = \sem{\gamma(t)}$. From $\Phi \models t \le s$ we get
$n \le \sem{\gamma(s)} \le m$, so the ledger has the step, and the new entry is
$\bud{a}{m - n}$. The result $\tok{k}$ has the type $\Bud{a}{\gamma(s \monus t)}$, because
$\sem{\gamma(s \monus t)} = \sem{\gamma(s)} - n \le m - n$.

\emph{Split.} As for \textsc{Spend}, $n \le m$, and the ledger has the step
$\evsplit(k, k', n)$ for a new name $k'$. Afterwards $k'$ has $n$ and $k$ has $m - n$, which
is at least $\sem{\gamma(s \monus t)}$. So the pair $(\tok{k'}, \tok{k})$ has the type
$\Bud{a}{\gamma(t)} \otimes \Bud{a}{\gamma(s \monus t)}$, and it owns
$\{k', k\} \subseteq O \cup N$.

\emph{Ask and Ask-Upto.} A grant enters $\ok{a}{x.\gamma(\varphi)}$ or
$\bud{a}{\sem{\gamma(t)}}$ under a new name $k$, and $\tok{k}$ has the type of the left
summand and owns $\{k\} \subseteq N$. A refusal returns $\inr()$.

\emph{Case, Let-P, Unpack, Let-B, If.} The operand has a value of the right shape by
Lemma~\ref{lem:values}(4), so the evaluation does not go wrong there, and the rest is as in
\textsc{Let-R}. In \textsc{If}, the branch that runs is the one whose added fact is true for
the numbers of $\gamma$. In \textsc{Unpack} the value is $\pack{n}{v_1}$ with $v_1$ of the type
$\gamma(\tau_1)[n/x]$. The environment $\gamma[x \mapsto n, y \mapsto v_1]$ still satisfies
$\Phi$, because $x \notin \Gamma$ and so no fact mentions $x$. The same holds in
\textsc{Let-B}. The type of the result does not mention $x$.

\emph{Iter.} By Lemma~\ref{lem:values}(4) the operand has a value of the type $\gamma(\tau)$
that owns a subset of $O$. Consider the loop $\iter{n}{v}{y}{\gamma(e)}$ for any ledger $L$
and any value $v$ with $L \vdash v : \gamma(\tau) \own O_v$. By induction on $n$ we show the statement of
the theorem for it, with $O_v$ in the place of $O$. For no turns nothing happens. Otherwise the body runs in the environment $\gamma[y \mapsto v]$, which owns $O_v$,
since the body has no other resource variable. The outer induction hypothesis gives a state
$v_1$ of the type $\gamma(\tau)$, and the inner one applies to the remaining turns from $v_1$.
The two compose as in \textsc{Let-R}.

\emph{Run.} The operand has a value $v$ of the closed type $\tau_1$ that owns
$O_v \subseteq O$. If the code is rejected, the result $\inr v$ has the type
$\tau_2 \oplus \tau_1$ and nothing changes. If it is accepted, the judgment
$\cdot; \cdot; y{:}\tau_1 \vdash e' : \tau_2$ has a derivation without \textsc{Run}, and
$L \vdash [y \mapsto v] : \cdot; \cdot; y{:}\tau_1 \own O_v$. The first part of the proof
applies to the code $e'$.

\emph{Sub.} By the induction hypothesis and Lemma~\ref{lem:values}(3).
\end{proof}

Items (2) and (3) are a frame property, as in separation logic~\citep{reynolds2002}. A part of
a program changes the ledger only at the names that its environment owns and at names that are
new, and its events mention no others. This is what lets a caller reason about a part that it
has not read.

\subsection{Tokens Can Be Erased}

\begin{lemma}[The ledger does not steer]
\label{lem:steer}
\begin{enumerate}
\item If $L, e \evto L', v, s$, then $L \lstep{s} L'$.
\item If $L, e \evto L', v, s$ and $\dom L_0 \subseteq \dom L$, then $L_0, e \evto \wrong$, or
  $L_0, e \evto L_0', v, s$ for some $L_0'$.
\item If $L, e \evto L', v, s$, then $\er{e} \evto^{\circ} \er{v}, \er{s}$. Conversely, if
  $\er{e} \evto^{\circ} u, r$, then for every ledger $L$ either $L, e \evto \wrong$, or
  $L, e \evto L', v, s$ for some $L'$, $v$, and $s$ with $\er{v} = u$ and $\er{s} = r$.
\end{enumerate}
\end{lemma}

\begin{proof}
Each by induction on the given derivation. A rule changes the ledger only by a step whose
event it adds to the trace, which gives (1). A rule consults the ledger only to ask whether a
step exists. No rule compares two names, and no value depends on an entry. For (2) and for the
second half of (3) we build the evaluation with the same numbers, the same answers, the same
code, and the same new names. A name that is new for $L$ is new for $L_0$, and the inclusion of
the domains is kept. Where the ledger at hand has the step that a rule needs, the rule
applies and gives the same value and event. Where it does not, the evaluation goes wrong.
\end{proof}

\begin{theorem}[Erasure]
\label{thm:erasure}
Let $\cdot; \cdot; \cdot \vdash e : \tau$. Then $e \evto^{\circ} u, r$ if and only if
$\emptyset, e \evto L, v, s$ for some $L$, $v$, and $s$ with $\er{v} = u$ and $\er{s} = r$.
\end{theorem}

\begin{proof}
From right to left by Lemma~\ref{lem:steer}(3). From left to right by the same lemma, since
by Theorem~\ref{thm:safety} no evaluation of $e$ from the empty ledger goes wrong.
\end{proof}

The ledger is state that exists for the proof, in the way of the ghost state of program
logics~\citep{jung2018}. An implementation keeps neither tokens nor a ledger. It may keep
them all the same, as a second line of defense, and then Theorem~\ref{thm:safety} says that the
second line is never reached.

\subsection{Every Action Has a Grant}

What does a trace of the ledger look like? The next definition says it without the ledger, in
terms that a person who reads the log can check.

\begin{definition}[Explained traces]
\label{def:explained}
In a trace, the event $\grant(k, E)$ \emph{introduces} the name $k$, and
$\evsplit(k, k', n)$ introduces the name $k'$ \emph{from} $k$. The names \emph{under} a name
$g$ are $g$ and every name that is introduced from a name under $g$. A trace is
\emph{explained} if every name is introduced at most once, and before every other event that
contains it, and
\begin{enumerate}
\item for every event $\grant(k, \ok{a}{x.\varphi})$ there is at most one action on $k$ and
  no split from $k$, and the action has the form $\evact(k, a, n)$ for a number $n$ that
  satisfies $\varphi$;
\item for every event $\grant(g, \bud{a}{m})$, every action on a name under $g$ is an action
  $a$, and the payloads of all these actions sum to at most $m$.
\end{enumerate}
\end{definition}

Every name is introduced by a grant or from a name that was introduced earlier, and by (1)
only budgets are split. So in an explained trace every action falls under exactly one grant,
and that grant covers it.

\begin{lemma}[The ledger explains]
\label{lem:ledger}
If $\emptyset \lstep{s} L$, then $s$ is explained.
\end{lemma}

\begin{proof}
By induction on the length of $s$ we show that $s$ is explained and that $L$ records what is
left: (a) the domain of $L$ is the set of the names that $s$ introduces; (b) a name that was
granted $\ok{a}{x.\varphi}$ has this entry if $s$ has no action on it, and $\gone$ otherwise;
(c) for a name $g$ that was granted $\bud{a}{m}$, every name $k$ under $g$ has an entry
$\bud{a}{m_k}$, and the sum of these $m_k$ and of the payloads of the actions on names under
$g$ is $m$. A grant or a split introduces a name that is not in the domain. By (b), an action
on an approval has a step only if the approval has not been used, and then the action and the
payload are as (1) asks. An action on a budget moves its payload from $m_k$ to the sum of the
payloads, and a split moves a part of $m_k$ to a new name under $g$; both keep (c). Since no
$m_k$ is negative, (c) gives (2).
\end{proof}

The lemma holds for every program, well typed or not: it is what the ledger enforces as a
monitor. For well-typed programs it holds without the monitor.

\begin{theorem}[Consent]
\label{thm:consent}
Let $\cdot; \cdot; \cdot \vdash e : \tau$. If $e \evto^{\circ} u, r$, then $r = \er{s}$ for a
trace $s$ that is explained.
\end{theorem}

\begin{proof}
By Theorem~\ref{thm:erasure} there is an evaluation $\emptyset, e \evto L, v, s$ with
$\er{s} = r$. By Lemma~\ref{lem:steer}(1), $\emptyset \lstep{s} L$, and by
Lemma~\ref{lem:ledger}, $s$ is explained.
\end{proof}

In words: take any execution of a well-typed program, with no ledger and no tokens. The events
of its log can be given names such that every action falls under a grant that the person made
before it. An approval pays for at most one action, of the kind that the person was shown and
with a payload that satisfies the predicate that the person was shown. The payloads under a
budget sum to at most its amount, however often it was split. The theorem quantifies over all
evaluations. It holds for a model that an attacker steers, for every reply of the world, and
for all the code that the model writes.

\subsection{A Part Does What Its Types Allow}

Theorem~\ref{thm:consent} is about whole programs. The next theorem is about a part: a branch,
the body of a loop, or code that the model has written. It bounds the actions of the part by
the \emph{types} of the tokens that the part is given.

\begin{definition}[Least ledger]
\label{def:least}
Let $L \vdash v : \tau \own O$. The \emph{least ledger} $\least{v}{\tau}$ has the domain $O$
and holds for every token what its type promises:
$\least{\tok{k}}{\Ok{a}{x.\varphi}} = [k \mapsto \ok{a}{x.\varphi}]$ and
$\least{\tok{k}}{\Bud{a}{t}} = [k \mapsto \bud{a}{\sem{t}}]$. For the other types it follows
the value: $\least{(v_1, v_2)}{\tau_1 \otimes \tau_2}$ is the union of $\least{v_1}{\tau_1}$
and $\least{v_2}{\tau_2}$; $\least{\inl v}{\tau_1 \oplus \tau_2} = \least{v}{\tau_1}$ and
likewise for $\inr$; $\least{\pack{n}{v}}{\exists x.\tau} = \least{v}{\tau[n/x]}$; and it is
empty for numbers and for $()$. For an environment, $\least{\gamma}{\Delta}$ is the union of
the ledgers $\least{\gamma(y)}{\gamma(\Delta(y))}$.
\end{definition}

A value has its type in its least ledger: $\least{v}{\tau} \vdash v : \tau \own O$, by
induction on Definition~\ref{def:values}.

\begin{theorem}[Authority]
\label{thm:authority}
Let $\Phi; \Gamma; \Delta \vdash e : \tau$ and $L \vdash \gamma : \Phi; \Gamma; \Delta \own O$,
and let $L, \gamma(e) \evto L', v, s$. Then the least ledger $\least{\gamma}{\Delta}$ has the
steps of $s$.
\end{theorem}

\begin{proof}
Let $L_0 = \least{\gamma}{\Delta}$. Its domain is $O$, which is a subset of $\dom L$, and
$L_0 \vdash \gamma : \Phi; \Gamma; \Delta \own O$. By Lemma~\ref{lem:steer}(2), an evaluation
of $\gamma(e)$ from $L_0$ goes wrong or ends with the same trace $s$. By
Theorem~\ref{thm:safety}, applied to $L_0$, it does not go wrong. By
Lemma~\ref{lem:steer}(1), $L_0 \lstep{s} L_0'$ for some $L_0'$.
\end{proof}

The theorem speaks about the ledger that the types describe and not about the ledger that
exists. The least ledger is built from the empty one by one grant for each of its names, so
Lemma~\ref{lem:ledger} applies to these grants followed by $s$. It gives three statements about
the trace $s$ of a part:
\begin{itemize}
\item An approval that the part holds at the type $\Ok{a}{x.\varphi}$ pays for at most one
  action, an $a$ with a payload that satisfies $\varphi$. This is the predicate of the type,
  also where the person agreed to more.
\item The actions under a budget that the part holds at the type $\Bud{a}{t}$, and under
  everything that is split from it, are actions $a$ whose payloads sum to at most $t$, also
  where more is left.
\item Every other action of the part falls under a grant that the person made while the part
  ran.
\end{itemize}

\begin{corollary}[Code that nobody has read]
\label{cor:code}
Let $L \vdash v : \tau_1 \own O$, let $y{:}\tau_1 \vdash_{\code} e : \tau_2$, and let
$L, e[v/y] \evto L', v', s$. Then $\least{v}{\tau_1}$ has the steps of $s$. In particular, code
that is accepted for $\Bud{a}{n} \multimap \tau_2$ performs, apart from what the person grants
it while it runs, only actions $a$ whose payloads sum to at most $n$.
\end{corollary}

By the proof of Theorem~\ref{thm:safety}, the hypothesis on $v$ holds at every \textsc{Run-Yes}
in an evaluation of a well-typed program. So a program that passes a budget at the type
$\Bud{\pay}{2}$ knows that the code will spend at most $2$ of it, before the code exists.

\section{Testing the Claims}
\label{sec:tests}

The artifact is one Python script without dependencies. It contains the checker, the two
semantics, Definition~\ref{def:explained} written as a function that does not use the ledger,
a generator of random programs, and the tests. This section says what was tested and what came
out. All numbers are those of one fixed seed, and \texttt{make -C artifact check} computes them
again and compares.

\subsection{Programs}

The generator builds programs of type $\Base$ by following the typing rules. It keeps the
three contexts and chooses among the forms that they allow: questions for approvals and
budgets, actions, tests, splits, sequences, loops that carry a budget of a changing amount or
an approval that may be used in a later turn, and calls of $\kw{run}$ with an approval or a
budget that is narrowed first. A payload is found by a search for a term for which the facts
entail the side condition. Programs have a nesting depth of at most five, and the median
program has \exMedianNodes{} nodes in its syntax tree.

A switch makes the generator \emph{slip}. With a probability of one in ten, at each place where
a rule constrains a choice, it ignores the constraint: it uses a token that is gone, picks a
payload without looking at the facts, names another action, hands code a type that the
argument does not have, lets the body of a loop use a token from outside, or binds a name
again.

The model, the world, and the person are a random stream. Numbers are uniform, and the person
agrees three times out of four. Code that the model writes comes from the same generator. Half
of the time it is written for the type that was asked. Otherwise it is written for a wider
type, such as a budget of $5$ where $2$ was given, or with slips.

\subsection{The Theorems}

We took $3{,}000$ programs from the generator without slips, and $3{,}000$ programs from the
generator with slips that the checker accepts. Of the latter, \exLucky{} contain a slip and are
accepted all the same, because the constraint that was ignored happens to hold. Each of the
\exPrograms{} programs ran three times with the ledger, and each execution was repeated
without tokens on the same random stream. The tests are these.
\begin{itemize}
\item Theorem~\ref{thm:safety}: the ledger refuses no step, and no value has a wrong shape.
\item Theorem~\ref{thm:erasure}: the execution without tokens returns the erasure of the value
  and the erasure of the trace.
\item Theorem~\ref{thm:consent}: the trace is explained.
\item The frame of Theorem~\ref{thm:safety}, wherever accepted code runs and at every turn of
  a loop: the tokens in the result are owned by the argument or new, no entry outside the
  names that the argument owns has changed, and no event names another token.
\item Theorem~\ref{thm:authority} and Corollary~\ref{cor:code}, at the same places: the least
  ledger of the argument at its type has the steps of the trace of the part.
\end{itemize}
In the \exRuns{} executions the person was asked \exAsks{} times and granted \exGrants{}
tokens, the programs performed \exActs{} actions, and the model wrote \exWritten{} pieces of
code, of which \exAccepted{} were accepted and ran. No test failed.

\subsection{Every Side Condition Is Needed}
\label{sec:needed}

\begin{table}[t]
\caption{Twelve checkers with one side condition removed. Each accepts a program that the
ledger refuses in one of three executions. The last column is the number of random programs
that were tried until the first such program.}
\label{tab:broken}
\small
\begin{tabular}{@{}llr@{}}
\toprule
Rule & What the broken checker allows & Programs tried \\
\midrule
division of $\Delta$ & a token is still there after it was used & \exTryContraction \\
\textsc{Case}, \textsc{If} & a token counts as used only if both branches use it & \exTryBranch \\
\textsc{Iter} & the body of a loop uses a token from outside & \exTryLoop \\
\textsc{Let-B}, \textsc{Unpack} & a base variable is bound a second time & \exTryShadow \\
\textsc{Act} & any payload under an approval & \exTryAct \\
\textsc{Act}, \textsc{Spend} & a token for one action serves for another & \exTryAction \\
\textsc{Spend} & any payload under a budget & \exTrySpend \\
\textsc{Split} & a part that is larger than the budget & \exTrySplit \\
\textsc{S-Ok} & an approval serves as any approval for the same action & \exTryAttenuate \\
\textsc{S-Bud} & a budget serves as any budget for the same action & \exTryAttenuateBudget \\
\textsc{If} & the second branch assumes the test and not its negation & \exTryIfFact \\
\textsc{Run-Yes} & code runs without the check & \exTryRun \\
\bottomrule
\end{tabular}
\end{table}

A type system can be sound because it rejects too much. To see that the conditions of
Figure~\ref{fig:typing} are not of that kind, we removed them one at a time.
Table~\ref{tab:broken} lists \exRules{} broken checkers. For each one the script draws programs
from the generator with slips, here with a probability of $0.15$, until it finds one that the
broken checker accepts and that the ledger refuses in one of three executions. It finds one for every broken checker. The slowest
case needs \exSlowest{} programs. It is the rule for branches, where the counterexample must
use a token in one branch only and use it again after the branches have joined.

The fourth row is the condition that a base variable is bound once. We had not planned it. The
first version of the checker lacked it, and the lack is a real hole: a test establishes a fact
about $x$, a new $x$ is bound, and the action under the fact uses the new one.

\subsection{A Ledger at Run Time Is Late}
\label{sec:late}

The ledger of Section~\ref{sec:ledger} can run as a monitor, and for programs that nobody
checks it is the only protection. How does it compare? We took $3{,}000$ programs from the
generator with slips that the checker rejects, and ran each of them twenty times with the
ledger. Programs that bind a name twice were left out, since that is a matter of the checker
and not of consent.

Every trace that the ledger let through was explained, as Lemma~\ref{lem:ledger} says. The
ledger stopped \exStopped{} of the \exRejectedRuns{} executions, which is \exStoppedPercent{}
percent. In \exLate{} of the executions that it stopped, \exLatePercent{} percent, at least one
action had been performed before the stop. The program had paid, sent, or deleted something,
and it could not go on. The checker rejects the same programs before they start. For code that
the model writes, that is the difference between a piece of code that is turned down, with its
argument handed back, and a task that is left half done.

The comparison has another side. The checker is conservative: \exQuiet{} of the $3{,}000$
programs, \exQuietPercent{} percent, were never stopped in twenty executions. Such a program
either contains a step that the ledger would refuse, on a path that twenty executions did not
take, or it is safe on every path for a reason that the rules do not see, such as two uses of
a token in branches that exclude each other. We did not separate the two cases.

\subsection{What the Tests Do Not Show}

The tests run programs of depth at most five over six numbers, and three executions of each.
They cannot find an error that needs a larger program, or a rare coincidence of answers. They
test the checker of the artifact, which accepts a subset of the judgments of
Figure~\ref{fig:typing}; rules that it never uses in their general form, such as \textsc{Sub}
at an arbitrary place, are covered by the proofs only. And the tests share the reading of the
definitions with the proofs: both were written by the same author.

\section{Limits}
\label{sec:limits}

\paragraph{The person.} The theorems begin at the answer. Whether the person read the
question, understood the predicate, or was worn down by the tenth dialog is outside the
calculus~\citep{felt2012}. Budgets reduce the number of questions, at the price that the
person no longer sees every payload. The calculus makes this price visible in a type, and it
does not say which price is right.

\paragraph{Payloads, not the world.} A predicate speaks about the payload of an action. It
does not speak about the world at the time of the action. An approval to delete the file of a
given name says nothing about what the file contains when it is deleted. The gap between check
and use is closed for what the program computes and stays open for what the world does in the
meantime.

\paragraph{Time.} A token does not expire, and the person cannot take it back. An expiry date
would be an index of the type together with a clock, and an action could then fail and would
have to return a sum. Revocation is dynamic by nature. Both give up the property that a
well-typed program stops only where it asks.

\paragraph{Questions.} Nothing bounds how often a program asks, and code that the model writes
may ask too. A question costs attention, and attention is the scarce resource here. A token
that permits a number of questions would fit the calculus, and we have not worked it out.

\paragraph{The language.} Programs are first order with bounded loops, code may not call
$\kw{run}$, and every program ends. For programs that run without end the theorems would
speak of the finite parts of a run; we expect them to hold and have not proved it. More
important in practice, code that the model writes has to be written in this calculus, while
agents write Python. A checker for a subset of a real language, with the same two token types,
is engineering that we have not attempted.

\paragraph{Trust.} What has to be trusted is the checker, the two rules that perform an
action, and the absence of any other way to act: if code can reach the network without going
through $\kw{act}$, no type helps. The proofs are on paper, and the artifact tests them and
does not prove them.

\section{Related Work}
\label{sec:related}

\paragraph{Capabilities.} A token is a capability in the sense of
\citet{dennis1966}: evidence of a permission that a program must present in order to act, and
that it cannot make. Capabilities serve the principle of least
privilege~\citep{saltzer1975} and remove the confused deputy~\citep{hardy1988}, and languages
built on them pass authority only by passing references~\citep{miller2006}. Macaroons let the
holder of a credential add caveats before passing it on~\citep{birgisson2014}; our subtyping
is this attenuation, and a split is attenuation for amounts. These are mechanisms of the run
time. Here the tokens are checked statically and erased, and copying is excluded by the types
and not by cryptography. That only \textsc{Ask} creates a token is the sealing
of~\citet{morris1973}, with the person as the only holder of the seal.

\paragraph{Linear types and resources.} Linear logic~\citep{girard1987} gave the idea of a
hypothesis that is used once, and affine types make it practical in a
language~\citep{tov2011}. Vault~\citep{deline2001} tracks keys linearly to enforce protocols
on resources, and the capability calculus~\citep{crary1999} and $\mathrm{L}^3$
\citep{morrisett2005} track capabilities for memory in types. Typestate~\citep{strom1986} lets the
type of a value change as it is used; an approval has two states and a budget has a counter.
Closest in spirit for budgets is Move, in which money is a linear resource and conservation is
a theorem~\citep{blackshear2020}. Our subject differs in what a token stands for, which is an
answer of a person together with what the person was shown, and in who writes the code that
holds it.

\paragraph{Quantities in types.} Indexed types in the style of \citet{xi1999} supply the
predicates and the amounts. Bounded linear logic counts the uses of a
hypothesis~\citep{girard1992}, and graded modal types generalise the count to a
semiring~\citep{orchard2019}. A budget is closer to the potential of \citet{hofmann2003}, which
pays for heap cells where ours pays for actions, and its splitting resembles that of fractional
permissions~\citep{boyland2003}.

\paragraph{Separation and ghost state.} The proof of Theorem~\ref{thm:safety} is the frame
rule of separation logic~\citep{reynolds2002} in small: owned sets are disjoint, and a part
changes only what it owns. The ledger is ghost state~\citep{jung2018}, and
Theorem~\ref{thm:erasure} is its erasure. The least ledger of Theorem~\ref{thm:authority} uses
the fact that a specification which holds for every ledger that fits the types holds for the
smallest one.

\paragraph{Monitors.} The policies that a monitor can enforce are the safety properties
\citep{schneider2000}, and being explained is one of them. Edit automata extend monitors with
the power to suppress and insert actions~\citep{ligatti2005}. \citet{walker2000} turns a security
automaton into the signature of a typed language, so that a host can check before it runs code
that the code respects the policy; our type system does this for the ledger. Our monitor's state is
made of tokens that the program holds, which is what makes delegation expressible: a part can
be given some of the state.

\paragraph{Authorization logics.} Logics of access control reason about who said what
\citep{abadi1993}, and proofs in such logics can be carried by requests~\citep{appel1999} and
by programs, as in AURA~\citep{jia2008}. \citet{garg2006} add linearity to obtain credentials
that are consumed by use. An approval is such a credential, issued only by the person and
specialised to a predicate on one payload; a budget is one that is consumed in parts.

\paragraph{Code from an untrusted author.} In proof-carrying code~\citep{necula1997} the
producer of code supplies a proof that the consumer checks. Here the model supplies code only,
and the type against which it is checked is the policy. This is an instance of generate and
check: the checker is sound, so an error of the model becomes a refusal that the program must
handle~\citep{tpp3}.

\paragraph{Agents.} Agents that act by writing code are common~\citep{wang2024}.
CaMeL~\citep{debenedetti2025} attaches capabilities to the values of such code in order to
control where untrusted data may flow, and a companion paper derives from the origin of a
payload whether an action needs a question at all~\citep{tpp7}. Those systems answer the
question of who chose a payload. This paper answers the question of whether the person agreed
to it, and what the agreement covers once it has been given.

\section{Conclusion}

An answer of a person is a small thing that is easy to stretch: to another payload, to a second
use, to a part of the agent that was not meant. We have treated it as a value that cannot be
copied, with a type that says what it covers. The type system that results is small, two
token types and a handful of side conditions, and its theorems say what one would hope for. Every
action falls under a grant, what was shown is what is done, a budget holds however it is
split, and code that nobody has read can do no more than the type at which it was accepted.
None of this needs a monitor while the program runs. What it needs is that programs, and the
code that models write, are checked before they act.

\bibliographystyle{ACM-Reference-Format}
\bibliography{refs}

\end{document}
